\section{A mass gap for homogeneous dispersion}\label{sec:r8-homogeneous}

Fix an integer $N\ge1$, a real order $m>1$, and a real, even symbol
$a\in C^{N+2}(\mathbb R^N\setminus\{0\})$ such that
\begin{equation}\label{eq:r8-symbol}
 a(t\xi)=t^m a(\xi)\quad(t>0),\qquad
 c|\xi|^m\le a(\xi)\le C|\xi|^m
\end{equation}
for positive constants $c,C$. Set
\[
 s=\frac m2,\qquad r=\frac mN,\qquad
 \alpha=2r,\qquad p=2(r+1),
\]
and define
\[
 A_a(f)=\int_{\mathbb R^N}a(\xi)|\widehat f(\xi)|^2\,d\xi,\qquad
 M(f)=\|f\|_2^2,\qquad
 E_a(f)=\frac12 A_a(f)-\frac1p\|f\|_p^p.
\]
The Fourier transform is unitary. The critical mass $M_{c,a}>0$ is defined
by the sharp Gagliardo--Nirenberg inequality
\begin{equation}\label{eq:r8-gn}
 \|f\|_p^p\le\frac p2
 \left(\frac{M(f)}{M_{c,a}}\right)^r A_a(f),
 \qquad f\in H^{m/2}(\mathbb R^N;\mathbb C).
\end{equation}
The kinetic form is equivalent to the homogeneous $H^{m/2}$ norm.
In particular, Sobolev interpolation proves that the sharp constant in
\eqref{eq:r8-gn} is finite and positive.

\begin{theorem}[Mass gap for homogeneous dispersion]\label{thm:r8-homogeneous-gap}
There exists $\varepsilon_{N,m,a}>0$ such that every nonzero weak
$H^{m/2}$ solution of
\begin{equation}\label{eq:r8-profile}
 a(D)U+\omega U+i v\cdot\nabla U-|U|^{2m/N}U=0,
 \qquad (\omega,v)\in\mathbb R\times\mathbb R^N,
\end{equation}
with $E_a(U)<0$ satisfies
\[
 M(U)\ge M_{c,a}+\varepsilon_{N,m,a}.
\]
The gap is uniform in frequency and velocity. No radial symmetry,
finite spatial moment, optimizer uniqueness, or real-part
nondegeneracy assumption is required.
\end{theorem}

Thus the exclusion statement applies to all fractional Laplacians
$a(D)=(-\Delta)^{m/2}$ of order $m>1$, all higher integer orders, and
all real-even homogeneous anisotropies satisfying \eqref{eq:r8-symbol}.
It does not assert existence of negative-energy traveling waves for
every such operator. The sharp mass and the gap depend on the symbol.

\subsection{Regularity at the energy level}

\begin{lemma}[Profile regularity]\label{lem:r8-regularity}
Every weak $H^{m/2}$ profile in \eqref{eq:r8-profile} belongs to
$L^\infty\cap H^{m+1}$.
\end{lemma}
\begin{proof}
Since $m>1$, for the fixed velocity $v$ one can choose $\mu>0$ so that
\[
 a(\xi)-v\cdot\xi+\mu\asymp1+|\xi|^m.
\]
The resolvent
\[
 R_\mu=(a(D)+iv\cdot\nabla+\mu)^{-1}
\]
maps $L^q$ to the Bessel potential space $H^{m,q}$ for every
$1<q<\infty$. Indeed, the multiplier
\[
 \frac{(1+|\xi|^2)^{m/2}}
 {a(\xi)-v\cdot\xi+\mu}
\]
satisfies the Mihlin derivative bounds. The possible nonsmoothness
of $a$ at zero is harmless: homogeneity gives the required bounds
after multiplication of each derivative by the corresponding power
of $|\xi|$. The stated $C^{N+2}$ regularity away from zero suffices.
None of the regularity constants below must be uniform in $v$.

Write
\begin{equation}\label{eq:r8-resolvent}
 U=R_\mu(|U|^\alpha U)+(\mu-\omega)R_\mu U.
\end{equation}
We first obtain $U\in L^\infty$. If $N\le m$, the energy space
embeds into all finite $L^q$; it embeds directly into $L^\infty$
when $N<m$. In the endpoint case take a sufficiently large finite $q$.
If $N>m$, start with
\[
 q_0=\frac{2N}{N-m},\qquad
 \beta_0=\frac1{q_0}=\frac12-\frac{m}{2N}.
\]
For $U\in L^{q_j}$ the nonlinear term is in
$L^{q_j/(1+\alpha)}$. Sobolev embedding in
\eqref{eq:r8-resolvent} improves reciprocal integrability according to
\[
 \beta_{j+1}=(1+2m/N)\beta_j-m/N
\]
as long as the right side is positive. The linear resolvent term has
the better reciprocal output exponent $\beta_j-m/N$ and also remains
in $L^2$; interpolation therefore gives the desired common output
exponent. The ideal iterates are
\[
 \beta_j=\frac12-\frac{m}{2N}(1+2m/N)^j.
\]
All nonlinear input exponents exceed one, since
$q_0>1+2m/N$ and the $q_j$ increase. Intermediate exponents can be
weakened slightly to avoid an endpoint. After finitely many steps,
$q_j/(1+\alpha)>N/m$, so the resolvent and Sobolev embedding give
$U\in L^\infty$.

Now $|U|^\alpha U\in L^2$, and \eqref{eq:r8-resolvent} gives
$U\in H^m\subset H^1$. The map
$F(z)=|z|^\alpha z$ on $\mathbb C\simeq\mathbb R^2$ is $C^1$
with bounded derivative on bounded sets. The Sobolev chain rule
therefore gives $F(U)\in H^1$. A final application of the resolvent
yields $U\in H^{m+1}$.
\end{proof}

\subsection{A nonlocal dilation identity without spatial moments}

Let $\chi\in C_c^\infty(\mathbb R^N;\mathbb R)$ equal one near zero, and
put
\[
 f_R(x)=x\chi(x/R)=R f(x/R),\qquad
 L_R=f_R\cdot\nabla+\frac12\operatorname{div}f_R,\qquad R\ge1.
\]
The operator $L_R$ is skew-symmetric on $L^2$.

\begin{lemma}[Cutoff commutator]\label{lem:r8-commutator}
The commutators $[a(D),L_R]$ are uniformly bounded from
$H^{m/2}$ to $H^{-m/2}$, and for every $u\in H^{m/2}$,
\begin{equation}\label{eq:r8-commutator-limit}
 \langle u,[a(D),L_R]u\rangle\longrightarrow mA_a(u)
 \qquad(R\to\infty).
\end{equation}
\end{lemma}
\begin{proof}
Apart from the fixed Fourier normalization factor, the kernel is
\[
 \frac i2\bigl(a(\xi)-a(\eta)\bigr)
 (\xi+\eta)\cdot\widehat f_R(\xi-\eta).
\]
Homogeneity and $C^1$ regularity on the sphere give
\[
 |a(\xi)-a(\eta)|
 \le C|\xi-\eta|(|\xi|+|\eta|)^{m-1}.
\]
For $k=\xi-\eta$, it follows that
\[
 \frac{|a(\xi)-a(\eta)|\,|\xi+\eta|}
 {\langle\xi\rangle^{m/2}\langle\eta\rangle^{m/2}}
 \le C|k|\langle k\rangle^{m/2}.
\]
Here $(|\xi|+|\eta|)^m$ is bounded by a constant times
$\langle\xi\rangle^m+\langle\eta\rangle^m$, and
$\langle\xi\rangle/\langle\eta\rangle$ and its reciprocal are bounded
by a constant times $\langle k\rangle$.
Since $\widehat f_R(k)=R^{N+1}\widehat f(Rk)$,
\[
 \int |k|\langle k\rangle^{m/2}|\widehat f_R(k)|\,dk
 =
 \int |\ell|\langle\ell/R\rangle^{m/2}|\widehat f(\ell)|\,d\ell
 \le C_f.
\]
Young's inequality proves the asserted uniform form bound.

For Schwartz $u$, $L_Ru\to Lu$ in $L^2$, where
$L=x\cdot\nabla+N/2$. Therefore
\[
 \langle u,[a(D),L_R]u\rangle
 =2\operatorname{Re}\langle a(D)u,L_Ru\rangle
 \longrightarrow
 2\operatorname{Re}\langle a(D)u,Lu\rangle=mA_a(u).
\]
The last equality follows from Fourier integration by parts and
$\xi\cdot\nabla a(\xi)=ma(\xi)$. The uniform form bound and
Schwartz density extend this limit to every $u\in H^{m/2}$.
\end{proof}

The same kernel argument, applied to the degree-one symbol
$-v\cdot\xi$, gives
\begin{equation}\label{eq:r8-drift-limit}
 \langle u,[iv\cdot\nabla,L_R]u\rangle
 \longrightarrow-v\cdot P(u).
\end{equation}
For this assertion the form estimate uses $H^{1/2}$, which contains
the energy space continuously because $m>1$.

\begin{lemma}[Profile identities]\label{lem:r8-identities}
Every profile in \eqref{eq:r8-profile} satisfies
\begin{align}
 v\cdot P(U)&=2mE_a(U),\label{eq:r8-pohozaev}\\
 \omega M(U)&=rA_a(U)+(2m-p)E_a(U).\label{eq:r8-frequency}
\end{align}
Writing $U=A+iB$ with real $A,B$, one also has
\begin{equation}\label{eq:r8-phase-energy}
 E_a(U)=\frac1m
 \left\langle B,\bigl(a(D)+\omega-|U|^{2r}\bigr)B\right\rangle.
\end{equation}
\end{lemma}
\begin{proof}
Lemma~\ref{lem:r8-regularity} makes $L_RU\in H^{m/2}$ an admissible
test in the weak equation. The kinetic and drift pairings converge
by \eqref{eq:r8-commutator-limit} and \eqref{eq:r8-drift-limit}.
The frequency pairing vanishes. For the nonlinearity, integration
by parts gives
\[
 \operatorname{Re}\int |U|^\alpha U\,\overline{L_RU}
 =
 \left(\frac12-\frac1p\right)
 \int\operatorname{div}f_R\,|U|^p
 \longrightarrow\frac mp\|U\|_p^p.
\]
This integration is justified by $|U|^\alpha U\in L^2$,
$\nabla U\in L^2$, and the compact support of $f_R$.
The coefficients $\operatorname{div}f_R$ are uniformly bounded
and converge to $N$. The resulting identity is
\[
 \frac m2 A_a(U)-\frac12v\cdot P(U)-\frac mp\|U\|_p^p=0,
\]
which proves \eqref{eq:r8-pohozaev}. Pairing the profile equation
with $U$ gives
$A_a(U)+\omega M(U)-v\cdot P(U)-\|U\|_p^p=0$,
and then \eqref{eq:r8-frequency}.

Because the symbol is real and even, $a(D)$ preserves real
functions. The imaginary profile equation is
\[
 \bigl(a(D)+\omega-|U|^{2r}\bigr)B=-v\cdot\nabla A.
\]
The identity $P(U)=-2\int B\nabla A$ and
\eqref{eq:r8-pohozaev} yield \eqref{eq:r8-phase-energy}.
\end{proof}

\subsection{Compactness and phase coercivity}

\begin{lemma}[Compactness of normalized near-optimizers]
\label{lem:r8-compactness}
The sharp inequality \eqref{eq:r8-gn} has an optimizer.
If $M(f_n)\to M_{c,a}$, $A_a(f_n)=1$, and $E_a(f_n)<0$,
then after translations and passage to a subsequence,
$f_n$ converges strongly in $H^{m/2}$ to a normalized optimizer.
\end{lemma}
\begin{proof}
The sharp inequality implies
$\|f_n\|_p^p\to p/2$ and $E_a(f_n)\to0$.
Local Sobolev nonvanishing gives translations for which a subsequence
has a nonzero weak $H^{m/2}$ limit $Q$. Local Rellich compactness
also gives almost-everywhere convergence after a subsequence.
Write
\[
 m_1=M(Q)>0,\quad a_1=A_a(Q)>0,\quad
 m_2=M_{c,a}-m_1,\quad a_2=1-a_1.
\]
Hilbert norm splitting, the Brezis--Lieb lemma, and
\eqref{eq:r8-gn} give
\[
 M_{c,a}^r\le m_1^ra_1+m_2^ra_2,\qquad a_1+a_2=1.
\]
If $0<m_1<M_{c,a}$, both coefficients $m_j^r$ are strictly below
$M_{c,a}^r$, which is impossible. Hence $m_1=M_{c,a}$, and the
same inequality forces $a_1=1$. Strong energy-space convergence
and equality in the sharp inequality follow.
To prove attainment, normalize an arbitrary maximizing sequence
by amplitude and dilation so that $M(f_n)=M_{c,a}$ and
$A_a(f_n)=1$. Its nonlinear norms converge to $p/2$, and the
identical splitting argument applies; negative energy is not
needed once that convergence is known.
\end{proof}

\begin{lemma}[Constant phase of optimizers]\label{lem:r8-constant-phase}
Every complex optimizer is a constant phase times a real function.
\end{lemma}
\begin{proof}
Write $f=A+iB$, where $A,B$ are real. Evenness of $a$ gives
$A_a(f)=A_a(A)+A_a(B)$. With
$C=(p/2)M_{c,a}^{-r}$, Minkowski, the sharp inequality on each real
component, and two-term H\"older yield
\begin{align*}
 \|f\|_p^2
 &\le\|A\|_p^2+\|B\|_p^2\\
 &\le C^{2/p}\left[
 A_a(A)^{2/p}M(A)^{1-2/p}
 +A_a(B)^{2/p}M(B)^{1-2/p}\right]\\
 &\le C^{2/p}A_a(f)^{2/p}M(f)^{1-2/p}.
\end{align*}
For an optimizer all inequalities are equalities. Strict convexity
of $L^{p/2}$ implies $B^2=cA^2$ almost everywhere, unless one
component vanishes. Repeat this equality argument after multiplying
$f$ by $e^{i\pi/4}$. If both signs of $B/A=\pm\sqrt c$ occurred
on sets of positive measure, the new squared components could
not be proportional. When $c=1$ the new components instead have
disjoint nonzero supports, again giving strict Minkowski inequality.
Thus only one sign is possible, and $f$ has constant phase.
\end{proof}

Let $Q$ be a real optimizer normalized by $M(Q)=M_{c,a}$ and
$A_a(Q)=1$. Its Euler--Lagrange equation is
\[
 a(D)Q+\omega_QQ-|Q|^{2r}Q=0,\qquad
 \omega_Q=r/M_{c,a}>0.
\]

\begin{lemma}[Strict phase coercivity]\label{lem:r8-phase-coercivity}
The real quadratic form
\[
 L_-=a(D)+\omega_Q-|Q|^{2r}
\]
is nonnegative, has kernel $\operatorname{span}_{\mathbb R}\{Q\}$,
and obeys
\[
 \langle b,L_-b\rangle\ge c_Q\|b\|_{H^{m/2}}^2
 \quad\text{if }b\text{ is real and }\langle b,Q\rangle=0,
\]
for some $c_Q>0$.
\end{lemma}
\begin{proof}
Retain $A=A_a(Q)$ and $M=M(Q)$ in the algebra, so that
$\omega_Q=rA/M$. For real $b\in H^{m/2}$ set
$B=M(b)$ and $C_1=A_a(b)$.
The sharp inequality applied to $Q+i\sqrt t\,b$ gives
\begin{equation}\label{eq:r8-kernel-comparison}
 F(t):=\int(Q^2+tb^2)^{r+1}
 \le(r+1)(1+tB/M)^r(A+tC_1),\qquad t\ge0.
\end{equation}
The right derivatives at zero prove $L_-\ge0$.
If $b$ belongs to its form kernel, the first derivatives agree, and
\[
 \int |Q|^{2r}b^2=C_1+rAB/M.
\]
The pointwise second-order remainder for $F$ is nonnegative.
Comparison of second-order remainders in
\eqref{eq:r8-kernel-comparison} and Fatou's lemma yield
\[
 D:=\int |Q|^{2r-2}b^4
 \le(r-1)AB^2/M^2+2BC_1/M.
\]
On the zero set of $Q$, the integrand is interpreted through
the remainder: it is infinite if $r<1$ and $b\ne0$, equals $b^4$
if $r=1$, and equals zero if $r>1$.
Cauchy--Schwarz gives
\[
 D\ge\frac{(C_1+rAB/M)^2}{(r+1)A}.
\]
Combining the inequalities yields
$(C_1-AB/M)^2\le0$. Consequently all inequalities are equalities.
Equality in Cauchy--Schwarz implies
$b^2=(B/M)Q^2$ on $\{Q\ne0\}$. Its integral already equals $B$,
so $b=0$ on the remaining zero set as well.
The form kernel is a linear space containing $Q$. Apply the same
conclusion to $b+Q$ and subtract the identity for $b^2$.
It follows that $b$ is one constant multiple of $Q$.

Finally, multiplication by $|Q|^{2r}$ is compact from $H^{m/2}$
to $H^{-m/2}$. Indeed $|Q|^{2r}\in L^{p/(2r)}$,
H\"older with two $L^p$ factors makes the tail small, and the
embedding into local $L^p$ is compact because $p$ lies strictly
below the critical Sobolev exponent. Thus $L_-$ is a compact
form perturbation of the coercive form $a(D)+\omega_Q$.
Nonnegativity and its one-dimensional kernel give the asserted
coercivity on the $L^2$ orthogonal complement of $Q$.
\end{proof}

\subsection{The mass-gap contradiction}

\begin{proof}[Proof of Theorem~\ref{thm:r8-homogeneous-gap}]
If no uniform gap existed, there would be profiles $U_n$ with
$E_a(U_n)<0$ and $M(U_n)\downarrow M_{c,a}$.
The scaling
\[
 U(x)\mapsto\lambda^{N/2}U(\lambda x),\qquad
 (\omega,v)\mapsto(\lambda^m\omega,\lambda^{m-1}v)
\]
preserves mass and energy sign, so normalize $A_a(U_n)=1$.
Lemmas~\ref{lem:r8-compactness} and \ref{lem:r8-constant-phase}
give, after translations and phases,
$U_n\to Q$ strongly in $H^{m/2}$, with $Q$ a real normalized
optimizer. Identity~\eqref{eq:r8-frequency} gives
$\omega_n\to r/M_{c,a}=\omega_Q$.
A further phase tending to zero ensures
$\langle\operatorname{Im}U_n,Q\rangle=0$.

Strong $L^p$ convergence implies
$|U_n|^{2r}\to|Q|^{2r}$ in $L^{p/(2r)}$.
By H\"older with two $L^p$ factors,
\[
 a(D)+\omega_n-|U_n|^{2r}\longrightarrow L_-
 \quad\text{in form norm }H^{m/2}\to H^{-m/2}.
\]
Write $B_n=\operatorname{Im}U_n$. Phase coercivity and
\eqref{eq:r8-phase-energy} imply, for large $n$,
\[
 E_a(U_n)
 =\frac1m\langle B_n,(a(D)+\omega_n-|U_n|^{2r})B_n\rangle
 \ge\frac{c_Q}{2m}\|B_n\|_{H^{m/2}}^2\ge0,
\]
a contradiction. No uniform bound on $v_n$ is needed: the exact
profile identities eliminate velocity from this final argument.
\end{proof}

\subsection{Local robustness in the uniform topology}

\begin{theorem}[Local uniform robustness of the gap]
\label{thm:r8-symbol-robustness}
Fix a symbol $a_0$ satisfying \eqref{eq:r8-symbol}.
There exist $\delta\in(0,1/2)$ and $\varepsilon>0$ such that, for every
real-even homogeneous symbol $a$ of degree $m$, individually
$C^{N+2}$ away from zero and satisfying
\[
 \sup_{\xi\in S^{N-1}}\left|\frac{a(\xi)}{a_0(\xi)}-1\right|
 <\delta,
\]
every negative-energy traveling profile for $a$ satisfies
\[
 M(U)\ge M_{c,a}+\varepsilon.
\]
No uniform bounds on derivatives of $a$, or compactness of this
symbol neighborhood, are required.
\end{theorem}
\begin{proof}
Suppose the assertion fails. Taking the proposed neighborhood size
and gap both to be $1/n$, choose symbols $a_n$ and
negative-energy profiles $U_n$ such that
\[
 \sup_{S^{N-1}}|a_n/a_0-1|<1/n,\qquad
 0<M(U_n)-M_{c,a_n}<1/n.
\]
Normalize $A_{a_n}(U_n)=1$ by critical scaling.
Homogeneity gives
\[
 (1-1/n)A_{a_0}\le A_{a_n}
 \le(1+1/n)A_{a_0}.
\]
Taking sharp constants in \eqref{eq:r8-gn} shows
$M_{c,a_n}\to M_{c,a_0}$.
The sharp inequality and negative energy also give
\[
 \|U_n\|_p^p\to p/2,\qquad E_{a_n}(U_n)\to0,\qquad
 A_{a_0}(U_n)\to1.
\]
The proof of Lemma~\ref{lem:r8-compactness}, with these converging
normalizations, applies to the limiting symbol. Hence, after
translations and phases, $U_n\to Q$ strongly in $H^{m/2}$,
where $Q$ is a real normalized optimizer for $a_0$.
The frequency identity gives $\omega_n\to r/M_{c,a_0}$.

The kinetic forms converge in norm from $H^{m/2}$ to $H^{-m/2}$
by uniform convergence of the symbols on the sphere.
Together with convergence of the nonlinear potentials and
frequencies, this transfers the strict phase coercivity for $Q$
to the sequence. After arranging
$\langle\operatorname{Im}U_n,Q\rangle=0$, the energy identity
\eqref{eq:r8-phase-energy} yields the same contradiction as before.
Profile regularity and the cutoff identity were established separately
for each symbol and each velocity; their constants never entered this
compactness argument. Uniform derivative bounds are therefore
unnecessary.
\end{proof}

\begin{proposition}[Compact families of symbols]\label{thm:r8-symbol-family}
Fix $N\ge1$ and $m>1$. Let $\mathcal K$ be a family of symbols
as in \eqref{eq:r8-symbol}, each individually real, even, and
$C^{N+2}$ away from zero, whose restrictions to $S^{N-1}$ form
a compact subset of $C^0(S^{N-1})$.
There is a single $\varepsilon_{N,m,\mathcal K}>0$ such that every
negative-energy traveling profile for every $a\in\mathcal K$
satisfies $M(U)\ge M_{c,a}+\varepsilon_{N,m,\mathcal K}$.
\end{proposition}
\begin{proof}
The neighborhoods supplied by
Theorem~\ref{thm:r8-symbol-robustness} form an open cover of
$\mathcal K$ in the uniform topology. Choose a finite subcover
and take the minimum of its positive gap constants.
\end{proof}

\begin{remark}
The mechanism is a restriction on coherent negative-energy motion:
near the sharp mass threshold, variational compactness forces a
profile close to a real optimizer, whose only null phase direction is
the constant phase. Removing that direction makes the phase form
strictly positive, whereas a negative-energy traveling profile requires
a negative value of the same form. This argument is independent of
radiality and does not use a maximum principle.
\end{remark}
